
\documentclass[a4paper,11pt]{article} 

\usepackage[spanish]{babel}
\usepackage[utf8]{inputenc}

\usepackage{geometry, amsmath, amssymb}
\geometry{margin=2.5cm}

\usepackage{hyperref}
\usepackage{microtype}
\usepackage{parskip}
\usepackage{fancyhdr}
\usepackage{graphicx}

% me parece que con esto manejo los colores
\usepackage{color, soul}

\pagestyle{fancy}
\fancyhf{}
\lhead{Guía de Ejercicios 5}
\rhead{Tecnología Digital VI}
\cfoot{\thepage}

\title{\vspace{-2cm}\textbf{Guía de Ejercicios 4}\\[0.2cm]
\large Tecnología Digital VI}
\author{Juan Ignacio Elosegui}
\date{Universidad Torcuato Di Tella \\ Junio 2026}

\begin{document}
\maketitle
\section*{Fundamentos de Redes Convolucionales (CNNs)}
\subsection*{Ejercicio 1}
\subsection*{Ejercicio 2}
\subsection*{Ejercicio 3}
\subsection*{Ejercicio 4}

\section*{Arquitecturas famosas y Transfer Learning}
\subsection*{Ejercicio 5}
\subsection*{Ejercicio 6}
\subsection*{Ejercicio 7}

\section*{Detección de objetos y segmentación (I2I)}
\subsection*{Ejercicio 8}
\subsection*{Ejercicio 9}
\subsection*{Ejercicio 10}

\section*{Interpretabilidad y manipulación}
\subsection*{Ejercicio 11}
\subsection*{Ejercicio 12}
\subsection*{Ejercicio 13}
\end{document}
